\chapter{EXPONENTIALS AND LOGARITHMS.}

The history of the multiplicative group \(R_{+}^{*}\) of real numbers \textgreater0 is closely linked to that of the development of the notion of powers of a number \textgreater0, and of the notations employed to indicate them. The conception of a ``geometric progression'' made up of the successive powers of the same number goes back to the Egyptians and the Babylonians; it was familiar to Greek mathematicians, and one finds already in Euclid (Elements, Book IX, prop. 11) a general statement equivalent to the rule \(a^{m}a^{n}=a^{m+n}\) for integer exponents \textgreater0. In the Middle Ages, the French mathematician N. Oresme (XIVth century) finds this rule again; it is also with him that appears for the first time the notion of fractional exponent \textgreater0, with a notation already close to ours and some rules for calculation (stated in general form) concerning them (for example the two rules that we now write \((ab)^{1/n}=a^{1/n}b^{1/n},(a^{m})^{p/q}=(a^{mp})^{1/q}\) {[}74{]}. But the ideas of Oresme were too much in advance of the Mathematics of his time to exercise an influence on his contemporaries and his treatise sunk rapidly into forgetfulness. A century later, N. Chuquet stated again Euclid's rule; he introduces as well an exponential notation for the powers of unknowns in his equations, and does not hesitate to make use of the exponent 0 and integer exponents \textless0.¹ This time (and although the work of Chuquet remained in manuscript and does not seem to have been circulated much), the idea of isomorphism between the ``arithmetic progression'' of the exponents, and the ``geometric progression'' of the powers, will no longer be lost sight of; extended to negative exponents and to fractional exponents by Stifel ({[}298{]}, fol.~35 and 249-250), it ends up finally with the definition of logarithms and the construction of the first tables, undertaken independently by the Scot J. Napier, in 1614-1620, and the Swiss J. Bürgi (whose work did not appear until 1620, although his conception went back to the first years of the XVIIth century). With Bürgi, the continuity of the isomorphism established between R and \(R_{+}^{*}\) is implicitly assumed by the use of interpolation in the handling of the tables; it is on the contrary explicitly formulated in the definition of

Napier (as explicitly at least as was allowed by the fairly vague concept of continuity that obtained at that time). \(^{2}\)

We do not need to dwell here on the services rendered by logarithms in the numerical Calculus; from the theoretical point of view, their importance dates especially from the beginnings of the infinitesimal Calculus, with the discovery of series expansions of \(\log(1+x)\) and of \(e^{x}\) , and of the differential properties of these functions (see pp.~172 ff.). As far as the definition of exponentials and logarithms is concerned, it will be limited, until the middle of the XIXth century, to assuming intuitively the possibility of extending by continuity to the set of real numbers the function \(a^{x}\) defined for all rational x; and it is only once the notion of real number is definitively made precise and deduced from that of rational number, that there will be thought of giving a rigorous justification to this extension.

